% Explainable production-adjusted electrical-energy baselines for industrial feed pelleting -- manuscript for Energy (Elsevier)
% Compile: pdflatex -> bibtex -> pdflatex -> pdflatex  (or latexmk -pdf main.tex)
% main_review.tex produces the Editorial Manager version (line numbers, double spacing).
\documentclass[preprint,12pt]{elsarticle}
\usepackage[T1]{fontenc}
\usepackage[utf8]{inputenc}
\usepackage{newtxtext,newtxmath}
\usepackage{amsmath,mathtools,graphicx,booktabs,tabularx,array}
\usepackage{microtype,xcolor,placeins}
\usepackage[hyphens]{url}
\usepackage[colorlinks=true,linkcolor=black,citecolor=black,urlcolor=blue!45!black,pdftitle={Explainable production-adjusted electrical-energy baselines for industrial feed pelleting},pdfauthor={Sousa Santos, Cabello Eras, Herrera Rodino, Arrabal-Campos}]{hyperref}
\usepackage{orcidlink}
\ifdefined\withlinenumbers\usepackage[mathlines]{lineno}\fi % set by main_review.tex
\urlstyle{same}
\setlength{\emergencystretch}{2em}
\renewcommand{\topfraction}{.9}\renewcommand{\bottomfraction}{.8}
\renewcommand{\textfraction}{.08}\renewcommand{\floatpagefraction}{.75}
\newcommand{\ind}{\mathbf{1}}
\DeclareMathOperator*{\argmin}{arg\,min}
\newcommand{\fig}[4]{\begin{figure}[!htbp]\centering\includegraphics[width=#3\linewidth]{figures/#1.pdf}\caption{#4}\label{fig:#2}\end{figure}}
\journal{Energy and AI}

\begin{document}
\ifdefined\withlinenumbers\linenumbers\fi
\begin{frontmatter}
\title{Explainable production-adjusted electrical-energy baselines for industrial feed pelleting}

% Author block (order as agreed in Overleaf). Affiliations in English; ORCID iDs from the public ORCID registry.
\author[cuc]{Vladimir Sousa Santos\orcidlink{0000-0001-8808-1914}}
\ead{vsousa1@cuc.edu.co}
\author[uco]{Juan Jos\'e Cabello Eras\orcidlink{0000-0003-0949-0862}}
\ead{juancabelloe@correo.unicordoba.edu.co}
\author[ucom]{Eder El\'{\i}as Herrera Rodi\~no}
\ead{eeherrera@correo.unicordoba.edu.co}
\author[ual]{Francisco Manuel Arrabal-Campos\orcidlink{0000-0002-5510-6297}\corref{cor1}}
\ead{fmarrabal@ual.es}
\cortext[cor1]{Corresponding author.}

\affiliation[cuc]{organization={Department of Energy, Universidad de la Costa},
  addressline={Calle 58 \# 55--66}, city={Barranquilla}, country={Colombia}}
\affiliation[uco]{organization={Department of Mechanical Engineering, Faculty of Engineering, Universidad de C\'ordoba},
  addressline={Carrera 6 No.\ 77--305}, city={Monter\'{\i}a}, country={Colombia}}
\affiliation[ucom]{organization={Master's Programme in Mechanical Engineering, Universidad de C\'ordoba},
  addressline={Carrera 6 No.\ 77--305}, city={Monter\'{\i}a}, country={Colombia}}
\affiliation[ual]{organization={Department of Engineering, Research Centre CIAIMBITAL, Escuela Superior de Ingenier\'{\i}a, University of Almer\'{\i}a},
  addressline={Carretera de Sacramento s/n, La Ca\~nada de San Urbano}, postcode={04120}, city={Almer\'{\i}a}, country={Spain}}


\begin{abstract}
This study asks whether an explainable, production-adjusted baseline can state how much electricity a feed-pelleting operation should have used, and whether operating histories, modern temporal models or pretrained foundation models add anything to it. Electricity use in pelleting varies with the quantity produced, the product, the line and the batch organization, so neither raw consumption nor a single specific energy consumption (SEC) figure says much about performance. We use 2,745 real records from three industrial lines (2,742 with positive energy) with dosed mass, dosing-batch count, product characteristics, line identity and calendar variables as inputs. Seven methods were compared on 800 chronologically ordered development targets; the selected Ridge regression was frozen, with its interval calibration, and evaluated once against line-specific SEC on 798 later records. Eight current-input, history-enabled, temporal and pretrained configurations were compared on 584 matched targets with 16 completed same-line operations as context. Ridge reduced test MAE from 61.29 to 52.45~kWh, 14.42\% below SEC, on all three lines; this is lower prediction error, not an energy saving. The Temporal Fusion Transformer reached 53.65~kWh against 53.55~kWh for current-input Ridge, with no evidence of paired improvement, and N-HiTS, TimesFM~2.5 and Chronos-2 were less accurate. Every estimate decomposes exactly into its inputs, and fixed 90\% and 95\% residual intervals covered 85.96\% and 93.23\% of the test records, with one line systematically underestimated. The baseline is ready for retrospective assessment and for a digital twin; prospective use needs verified inputs, line-specific acceptance criteria and a new consecutive period.
\end{abstract}

\begin{highlights}
\item Production-adjusted electricity baseline from 2,745 real feed-pelleting records
\item Frozen Ridge model cuts test MAE by 14.42\% relative to line-specific SEC
\item TFT, N-HiTS, TimesFM 2.5 and Chronos-2 add no accuracy on identical plant records
\item Exact additive decomposition traces each kWh estimate to recorded production inputs
\item Fixed 90\%/95\% residual intervals cover 85.96\% and 93.23\% of later test records
\end{highlights}

\begin{keyword}
feed pelleting \sep industrial energy baseline \sep production adjustment \sep chronological validation \sep prediction intervals \sep explainable Ridge regression
\end{keyword}
\end{frontmatter}

\section*{Nomenclature}
\begingroup\small
\noindent\begin{tabularx}{\linewidth}{@{}>{$}l<{$}X@{}}
\multicolumn{2}{@{}l}{\textit{Symbols}}\\
y_i,\ E_i & pelleting electrical energy of record $i$ (kWh) \\
\widehat y_i & point estimate of $y_i$ (kWh) \\
m_i & dosed mass of record $i$ (t) \\
b_i & number of dosing batches of record $i$ \\
\ell_i,\ c_i,\ p_i,\ s_i & pelleting line, product subfamily, presentation and bag/bulk category of record $i$ \\
h_i,\ d_i & origin hour (h) and weekday (0--6) of record $i$ \\
t_i & analytical origin of record $i$: earliest recorded stage start \\
a_j & recorded pelleting completion of record $j$ \\
\Delta_i & recorded pelleting duration of record $i$ (h) \\
P_i^{\mathrm{avg}} & duration-based average-power proxy of record $i$ (kW) \\
z_i & standardized and one-hot encoded input vector of record $i$ \\
\widehat b,\ \widehat\beta & fitted intercept and coefficient vector of the Ridge model \\
\widehat s_\ell & training-period specific energy consumption of line $\ell$ (kWh/t) \\
\mathcal T,\ \mathcal P & training set and partition interval \\
\mathcal H_i^{(16)},\ \mathcal H_i & 16-record operating history and recent-residual history of target $i$ \\
C_{ig} & contribution of variable group $g$ to the estimate of record $i$ (kWh) \\
\delta_i & shrunk recent-residual correction (kWh) \\
q_\alpha & calibrated half-width at miscoverage level $\alpha$ (kWh) \\
L_i,\ U_i & lower and upper interval limits of record $i$ (kWh) \\
r_i^{\mathrm{excess}},\ B_i & signed excess over the estimate (kWh) and out-of-interval flag \\
\end{tabularx}

\medskip
\noindent\begin{tabularx}{\linewidth}{@{}lX@{}}
\multicolumn{2}{@{}l}{\textit{Abbreviations}}\\
EnB, EnPI & energy baseline, energy performance indicator \\
IPMVP & International Performance Measurement and Verification Protocol \\
MAE, RMSE, WAPE & mean absolute error, root mean squared error, weighted absolute percentage error \\
IS & interval score \\
MLP & multilayer perceptron \\
SEC & specific energy consumption (kWh/t) \\
TFT & Temporal Fusion Transformer \\
N-HiTS & Neural Hierarchical Interpolation for Time Series \\
XReg & covariate-regression adapter of TimesFM \\
H1, H2, H3 & post-hoc extensions: line-dependent mass effect, product-dependent mass effect, recent-residual correction \\
GPU, CPU & graphics and central processing unit \\
\end{tabularx}
\endgroup

\section{Introduction}
Two pelleting operations can consume different amounts of electricity for entirely ordinary reasons: one processes more material, uses a different product formulation or runs on another machine. These differences matter when a plant tries to judge whether an operation consumed more energy than expected. A single kWh/t reference is convenient, but it can leave much of the variation in production conditions unexplained. An energy baseline should account for that variation while remaining clear enough for engineers to understand why its estimate changes.

The physical basis for such adjustment is well established. Feed properties and processing conditions jointly influence pellet quality and manufacturing behaviour \citep{thomas1996,thomas1997}. Conditioning changes the thermal and moisture state of the material before compaction \citep{thomas2020}, and both steam-conditioning rate and feed composition can affect the process and its mechanical energy demand \citep{skoch1981,bastiaansen2023}. Industrial measurement and verification must nevertheless distinguish changes in production from savings attributable to an intervention \citep{kissock2008}. Related work on automated energy baselines shows why predictive accuracy and uncertainty matter for that distinction \citep{granderson2016,gallagher2018}.

Whether operating histories improve on production-adjusted estimates remains an empirical question. Regularized regression and boosting can capture associations between production inputs and consumption \citep{hoerl1970,friedman2001}; temporal architectures such as TFT and N-HiTS can also use operating histories \citep{lim2021,challu2023}. Pretrained models, including TimesFM and Chronos, offer another route when local training data are limited \citep{das2024,ansari2024}. Their usefulness depends on what the industrial records actually represent. A sequence of completed operations, each with a different mass and duration, is not equivalent to a regularly sampled power signal. The prediction origin and the information available at that time must therefore be defined before comparing models \citep{hewamalage2023}.

Uncertainty is equally consequential. An estimate can be close on average and still miss particular machines or periods. Prediction intervals help express this variability, provided that coverage is considered together with interval width \citep{gneitingcal2007,gneiting2007}. Conformal calibration provides a practical starting point \citep{angelopoulos2021}, although temporal dependence and changing operating conditions complicate its guarantees \citep{barber2023}. Good coverage over a complete dataset also need not imply adequate coverage within each machine or product group \citep{barber2021}.

Here we use real records from three feed-pelleting lines to develop an explainable baseline for electrical energy per operation. We first compare production-adjusted models with a machine-specific energy-intensity reference and assess the selected model in a chronological test. We then ask whether completed same-machine histories improve estimation through linear, neural or pretrained temporal models. Finally, we examine the baseline's variable contributions and the scope for updating its point estimates and uncertainty using recent residuals. The aim is a reproducible statistical component that can later be incorporated into a digital twin, with a clear account of the evidence needed for prospective use. Figure~\ref{fig:framework} illustrates the selected baseline and shows how its output is read for an actual industrial record.

\fig{00_framework}{framework}{1}{The selected energy baseline. (a) Recorded production conditions are transformed into an electricity estimate; residuals from a later calibration period determine its uncertainty limits. (b) A real operation on Pellet~2 illustrates the output. The light segment is the fixed 95\% interval and the dark segment the 90\% interval; the circle marks the estimate and the diamond the observed energy. The example retains source row 2080 and uses its archived values. An excess above the point estimate is a model residual, not a measured saving or an adjudicated fault.}

\section{Materials and methods}
\subsection{Industrial records and model variables}
The source workbook contains 2,745 real industrial records and 32 named variables. Pellet~1, Pellet~2 and Pellet~3 identify machines, with 730, 1,148 and 867 records, respectively. Recorded stage starts extend from 2 January to 1 April 2025; the last pelleting completion is on 2 April. Each row aggregates an operation and includes dosed mass, batch counts, product classifications, stage timestamps, electricity, steam and operational summaries. There is no verified order identifier linking every row and stage to a unique physical batch. We therefore retain the source row as the observational unit and refer to it as a record.

The response $y_i$ is the field \texttt{KWH PELET}: pelleting electricity in kWh per record. It excludes the separately recorded electricity of milling and mixing and the thermal energy associated with steam. Dosed mass $m_i$ is obtained from \texttt{PESO DOS (kg)} by dividing by 1,000. The distinction between the response, apparent average power and specific energy is
\begin{equation}
 E_i=y_i\;[\mathrm{kWh}],\qquad
 P_i^{\mathrm{apparent}}=\frac{y_i}{\Delta_i}\;[\mathrm{kW}],\qquad
 \mathrm{SEC}_i=\frac{y_i}{m_i}\;[\mathrm{kWh/t}],
 \label{eq:units}
\end{equation}
where $\Delta_i$ would be a valid metering duration in hours. We learn $E_i$ and use SEC to construct a reference model. Apparent power is not treated as a verified measurement because the relationship between the meter and record time boundaries remains unresolved.

The historical origin $t_i$ is the earliest available start among milling, mixing and pelleting. This definition is retained throughout the analyses. Inputs comprise dosed mass, dosing batch count, machine, product subfamily, presentation, bag/bulk category and cyclic calendar coordinates (Table~\ref{tab:inputs}). With $b_i$ denoting batch count, $h_i$ the origin hour plus minutes divided by 60, and $d_i\in\{0,\ldots,6\}$ the weekday indexed from Monday, the numerical vector is
\begin{equation}
 x_i^{\mathrm{num}}=\left(m_i,b_i,\sin\frac{2\pi h_i}{24},\cos\frac{2\pi h_i}{24},
 \sin\frac{2\pi d_i}{7},\cos\frac{2\pi d_i}{7}\right).
 \label{eq:inputs}
\end{equation}
Four categorical fields complete the input set. In the frozen Ridge model, standardization of the six numerical coordinates and one-hot encoding of the categories produce 22 design columns. Scaling and category vocabularies are fitted on training data only; unknown categories receive an all-zero code for their field.

\begin{table}[!htbp]\centering\small
\caption{Complete dictionary of current model inputs. Original source labels are retained to disambiguate similarly named fields. Counts and ranges describe all 2,745 source rows.}\label{tab:inputs}
\begin{tabularx}{\linewidth}{@{}p{.20\linewidth}p{.24\linewidth}X@{}}\toprule
Input / units & Source and representation & Interpretation and recorded support \\\midrule
Dosed mass $m_i$ / t & \texttt{PESO DOS (kg)} divided by 1,000 & Material quantity entering the recorded operation; 2.69--36.78 t, median 14.98 t. It is the mass denominator used throughout. \\\addlinespace
Dosing batches $b_i$ / count & \texttt{N\_BACHES\_DOS}; numerical & One to ten, median five. Describes aggregation and production quantity; correlated with total mass. It is not a unique physical-cycle identifier. \\\addlinespace
Machine $\ell_i$ / category & \texttt{PELLT}; one-hot & Pellet 1 / 2 / 3; 730 / 1,148 / 867 records. The separate source field \texttt{LINEA} contains product families. \\\addlinespace
Product subfamily $c_i$ / category & \texttt{SUBLINEA}; one-hot & Nine labels spanning poultry, swine, cattle, equine and rabbit products. A product classification is not a measured formulation. \\\addlinespace
Presentation $p_i$ / category & \texttt{PRESENTACION}; one-hot & \texttt{PELETIZADO}: 2,065; \texttt{CROMBELIZADO}: 680. Recorded pelleted/crumbled presentation may also reflect product allocation. \\\addlinespace
Bag/bulk $s_i$ / category & \texttt{DES\_PRODUCTO\_TIPO}; one-hot & \texttt{ENSAQUE}: 2,420; \texttt{GRANEL}: 325. The analytical alias is ``salida\_programada'', but the values identify bag/bulk form, not a scheduled departure timestamp. \\\addlinespace
Origin hour / dimensionless pair & Earliest available stage start; sine and cosine & Two coordinates avoid an artificial discontinuity between 23:59 and 00:00. Seconds are not used in this encoding. \\\addlinespace
Origin weekday / dimensionless pair & Same origin; sine and cosine & Two coordinates describe recurring weekly associations; they do not measure operational state. \\\bottomrule
\end{tabularx}\end{table}


Mass captures the scale of the operation, while batch count describes how that mass is aggregated. Product subfamily and presentation can reflect differences in feed properties and processing requirements \citep{thomas1997,thomas2020,bastiaansen2023}. Calendar terms allow recurring associations with scheduling or product mix. Bag/bulk status is part of the common recorded feature set; its inclusion does not extend the target to packaging electricity. These variables describe production conditions, but their fitted effects need not be causal or independent of one another.

The workbook stores consolidated values without a capture-time or plan-version history. In particular, the total mass of a record may not have been known at its earliest stage start. The analysis consequently estimates consumption conditional on recorded production characteristics. Current-record durations, final production quantity, stage electricity, and complete steam, temperature and current summaries are excluded from the inputs because their availability at the origin is not established. For historical labels, pelleting completion provides the earliest plausible availability time; actual publication may occur later.

All 32 source fields were reconciled with the analytical table across all rows, yielding 87,840 passing cell comparisons under the documented numerical and whitespace tolerances. This establishes the fidelity of the copy, rather than validating the instrumentation. Three zero energy values are excluded from supervised targets and are not imputed. Fifteen auxiliary operational-summary values had been imputed during earlier preparation; none is used by the baseline. The source dictionary also preserves unresolved unit semantics: for example, the pressure field is labelled psi, without establishing gauge or absolute pressure. Steam mass alone is insufficient to derive thermal consumption.

The timestamps reveal substantial irregularity. Median pelleting durations are about 145, 96 and 123 minutes across the three machines. There are 351 within-machine interval overlaps: 75 on Pellet~1, 131 on Pellet~2 and 145 on Pellet~3. Milling ends after pelleting starts in 528 records, and mixing does so in 376; these counts may overlap. Such patterns warrant clarification of the aggregation boundaries, but are not sufficient grounds for deleting records as duplicates. The workbook contains neither dense power trajectories nor verified active/inactive labels. A separate supplied set of synthetic ten-second series is therefore kept outside the industrial evidence presented here.

\subsection{Models and chronological evaluation}
The original allocation separates development before 22 February, calibration from 22 February to 5 March, and testing from 6 March onward (Figure~\ref{fig:allocation}). Fitting and calibration require valid positive outcomes, valid durations, the required inputs, and completion before the next boundary. These rules retain 1,542 of 1,556 development candidates and 380 of 391 calibration candidates. All 798 test candidates are retained, including five operations lasting more than 24 hours.

\fig{01_records_and_partitions}{allocation}{1}{Recorded operations and the original chronological allocation. (a) Daily counts are shown separately for each machine on the same scale; background bands distinguish development, calibration and testing, and dashed lines mark 22 February and 6 March. Every source row is counted by its earliest recorded stage start. (b) Candidate and retained sample sizes account for the 25 exclusions before testing. The dates refer to origins; the final operation completes on 2 April. Days without records do not establish zero electricity use.}

Model development uses four expanding training blocks with validation weeks starting on 25 January, 1 February, 8 February and 15 February. The training sizes are 705, 924, 1,121 and 1,334, and the corresponding validation sizes are 208, 187, 205 and 200: 800 common targets over 28 observed origin dates. Training outcomes must finish before validation starts; the original validation rule also requires outcomes to finish before that week closes. This latter rule restricts long operations at the boundary. Chronological assessment follows the direction in which a model would be used \citep{tashman2000,cerqueira2020}, with the usual dependence assumptions and leakage concerns kept explicit \citep{bergmeir2018,kapoor2023}.

The simplest production-adjusted reference estimates one energy intensity per machine from training totals:
\begin{equation}
 \widehat{s}_\ell=\frac{\sum_{i\in\mathcal T:\ell_i=\ell}y_i}{\sum_{i\in\mathcal T:\ell_i=\ell}m_i},
 \qquad \widehat y_i^{\mathrm{SEC}}=m_i\widehat s_{\ell_i}.
 \label{eq:sec}
\end{equation}
We also include the training median energy per machine and a recent-SEC reference based on the median of the last five completed same-machine energy-to-mass ratios. The recent reference updates as outcomes become available during validation, whereas the static models keep their fitted parameters fixed.

Let $z_i$ contain the transformed current inputs. Ridge regression estimates an unpenalized intercept and regularized coefficients,
\begin{equation}
 (\widehat b,\widehat\beta)=\argmin_{b,\beta}
 \left\{\sum_{i\in\mathcal T}(y_i-b-z_i^\top\beta)^2+10\|\beta\|_2^2\right\},
 \qquad \widehat y_i=\max(0,\widehat b+z_i^\top\widehat\beta).
 \label{eq:ridge}
\end{equation}
The non-negative output respects the sign of energy consumption, although it does not impose an energy balance. Its additive form is useful for inspecting the estimate directly \citep{rudin2019}. Correlated inputs and uneven product allocation still limit the interpretation of individual coefficients, a difficulty shared with more general attribution methods \citep{aas2021,fisher2019}.

ExtraTrees, histogram gradient boosting and a multilayer perceptron use the same current inputs. Their configurations are fixed: 250 trees with minimum leaf size five for ExtraTrees; 200 boosting iterations, seven leaves, learning rate 0.05, minimum leaf size 20 and regularization 10; and two neural hidden layers of 32 and 16 units, Adam optimization, regularization 0.1 and at most 800 iterations. The neural predictions average seeds 11, 29 and 47; the tree methods use seed 2026. No equal-compute search is claimed for this original comparison. The selected Ridge model and its interval calibration were frozen before the first evaluation of the 798-record test.

We subsequently examined whether operating histories added useful information. For each current target, a temporal model receives 16 completed records from the same machine and partition interval. If $\mathcal P$ denotes that interval and $a_j$ the completion of historical record $j$, the context is
\begin{equation}
 \mathcal H_i^{(16)}=\operatorname{last}_{16,(a_j,j)}
 \{j:\ell_j=\ell_i,\ t_j\in\mathcal P,\ a_j<t_i,\ a_j<\sup\mathcal P\}.
 \label{eq:window}
\end{equation}
The operator orders by completion and then source-row identifier. Thus an earlier-starting but unfinished operation cannot enter the history. Contexts restart at each weekly validation boundary. Completed validation outcomes can inform later contexts within that week, but never update fitted weights. The same policy applies to inner validation. The resulting horizon is one nominated record, whose elapsed duration varies; no regular sampling grid is created.

Requiring 16 within-week donors removes 216 of the original 800 targets. All eight temporal and control variants are evaluated on the remaining 584: 149 from Pellet~1, 257 from Pellet~2 and 178 from Pellet~3, over 21 dates. Outer training sizes are 650, 869, 1,066 and 1,279, with validation sizes 151, 133, 151 and 149. Figure~\ref{fig:windows} gives the full design and an actual context. Restricting both targets and training eligibility makes the current-input Ridge control appropriate for this comparison; the original 800-record ranking remains a separate result.

\fig{07_windows}{windows}{1}{Temporal evaluation and one actual context. (a) Training expands over four blocks; green bars identify the outer validation weeks. Each training bar also reports the earlier training / preceding-week validation counts used for epoch selection. (b) The 16 donors for target row 823 on Pellet~1, ordered by completion. Segments span recorded pelleting start and finish, and dots mark completion. All donor outcomes precede the target origin on 28 January at 04:24. The target's 148.5 kWh is withheld until evaluation. Source rows remain identifiable throughout.}

Two Ridge controls separate the effect of history from model capacity. One uses current covariates; the other flattens those covariates, the 16 historical covariate vectors and their energy values, retaining penalty 10. TFT and N-HiTS receive the same underlying current and historical fields through their official NeuralForecast implementations \citep{nixtlatft,nixtlanhits}. TFT combines recurrent processing, variable selection and attention \citep{lim2021}; N-HiTS uses hierarchical interpolation blocks \citep{challu2023}. Here TFT has hidden dimension 32, four heads and dropout 0.1. N-HiTS has three blocks with two 64-unit dense layers each, pooling factors 2, 2 and 1, and interpolation downsampling factors of one. The one-record task does not exploit the full multi-horizon capabilities of either architecture.

Each neural family uses seeds 11, 29 and 47 and the same selection budget. An inner validation week immediately precedes the outer week. Optimization uses standardized-target MSE, Adam at $10^{-3}$, weight decay $10^{-4}$, batch size 64 and gradient clipping at one. Training runs for at most 40 epochs, with at least ten epochs examined and patience eight; the best inner-validation epoch may be earlier than ten. Each seed is then refitted on complete outer training for its selected epoch count. The reported output averages its three individually clipped predictions. The inner split fits its own transformations, and neither epochs nor seeds are chosen on outer validation.

\begin{table}[!htbp]\centering\small
\caption{Information supplied to the matched 584-record comparison. Current covariates mean the six numerical and four categorical fields in Table~\ref{tab:inputs}. All sequence contexts contain exactly 16 completed same-machine records.}\label{tab:models}
\begin{tabularx}{\linewidth}{@{}p{.25\linewidth}XX@{}}\toprule
Variant & Available information & Fitting / adaptation policy \\\midrule
Ridge: current & Current covariates & Supervised outer-training fit; penalty 10. \\
Ridge: full window & Current and historical covariates; historical energy & Same target eligibility; flattened window; penalty 10. \\
TFT; N-HiTS & Current and historical covariates; historical energy & One architecture per family; three seeds; inner-week epoch selection and outer-training refit. \\
TimesFM 2.5: energy only & Historical energy & Frozen pretrained weights; no local weight update. \\
TimesFM 2.5 + XReg & Current/historical covariates and historical energy & Native per-origin covariate regression, penalty 10, plus frozen TimesFM residual forecast. \\
Chronos-2: energy only & Historical energy & Frozen pretrained weights; median forecast. \\
Chronos-2 + covariates & Current/historical covariates and historical energy & Native covariate interface; median forecast; cross-query learning disabled. \\\bottomrule
\end{tabularx}\end{table}


TimesFM 2.5 uses the frozen 200M checkpoint, both alone and with its native XReg adapter \citep{das2024,timesfm25,timesfmrepo}. XReg fits a local covariate regression on the 16 historical records, adds the foundation model's residual forecast and uses penalty 10. Each origin is processed separately so that a regression cannot pool labels from another query whose outcome lies in its future. Chronos-2 is likewise evaluated with energy-only and native covariate inputs, using its median forecast and disabling cross-query learning \citep{ansari2025,chronosrepo}. Native categorical encoding is based on historical context, with the current target withheld. The current record occupies the interface's future-covariate position; this is an indexing convention rather than evidence that its consolidated plant inputs were known in advance.

The neural computations use float32 on an NVIDIA RTX PRO 5000 Blackwell GPU with 48 GB, with TF32 disabled. Ridge and the native JAX XReg solver run on CPU. TFT and N-HiTS were trained afresh, while foundation weights remained frozen. The preceding CPU run is retained for numerical comparison (Appendix~\ref{sec:computing}). These fixed local configurations test a specific representation of the industrial records; they do not establish the general ranking of model families on broader forecasting tasks \citep{aksu2024,zeng2023}.

We also tested three simpler extensions motivated by the original results. H1 allows the mass association to vary by machine, and H2 by product subfamily. With grouping variable $g_i$, their untruncated predictor is
\begin{equation}
 f_i=b+z_i^\top\beta+z_{m,i}\sum_{c\in\mathcal C}\gamma_c\ind(g_i=c),
 \label{eq:interaction}
\end{equation}
using the same Ridge penalty and training-only transformations. H3 instead corrects the fixed Ridge prediction using recent same-machine residuals. Let $e_j=y_j-\widehat y_j^{\mathrm{Ridge}}$, and let $\mathcal H_i$ contain up to 20 validation records completed before $t_i$, with the history reset weekly. For $n_i=|\mathcal H_i|$,
\begin{equation}
 \delta_i=\begin{cases}
 \displaystyle\frac{n_i}{n_i+10}\frac{1}{n_i}\sum_{j\in\mathcal H_i}e_j,&n_i\ge5,\\
 0,&n_i<5,
 \end{cases}
 \qquad \widehat y_i^{\mathrm{H3}}=\max(0,\widehat y_i^{\mathrm{Ridge}}+\delta_i).
 \label{eq:bias}
\end{equation}
The correction is shrunk when few outcomes are available and always uses residuals of the fixed model. Like the subsequent temporal benchmark, H1--H3 are exploratory development analyses. They do not replace the model evaluated in the original final test.

\subsection{Uncertainty and performance assessment}
The frozen Ridge model is calibrated using absolute errors from the 380 subsequent records. For miscoverage level $\alpha$, the empirical order statistic and prediction interval are
\begin{equation}
 q_\alpha=r_{(\lceil(n_{\mathrm{cal}}+1)(1-\alpha)\rceil)},\quad
 r_j=|y_j-\widehat y_j|,\quad
 I_i^{1-\alpha}=[\max(0,\widehat y_i-q_\alpha),\widehat y_i+q_\alpha].
 \label{eq:interval}
\end{equation}
The 90\% and 95\% half-widths are 93.0132 and 132.9576 kWh. We use this split-conformal-style construction as an empirical calibration method without assuming exchangeability of the industrial sequence.

A post-hoc extension updates the quantile using the latest 200 available absolute residuals, ordered by completion and source row. The pool starts with calibration residuals and then admits completed test outcomes, globally across machines; the point estimates remain fixed. Every donor must satisfy
\begin{equation}
 a_j<t_i,\qquad a_j=\text{recorded pelleting completion of donor }j.
 \label{eq:arrival}
\end{equation}
The actual publication time would replace this proxy in deployment. All 12,981 donor relationships for H3 and 159,600 for the interval update pass the recorded chronology check. The rolling quantile is a simple comparator, distinct from dedicated adaptive conformal algorithms \citep{xu2021,zaffran2022,gibbs2024}; the original test is already used and cannot provide an independent validation of a change motivated by its results.

Point performance is summarized by
\begin{align}
 \mathrm{MAE}&=N^{-1}\sum_i|\widehat y_i-y_i|,&
 \mathrm{RMSE}&=\sqrt{N^{-1}\sum_i(\widehat y_i-y_i)^2},\\
 \mathrm{Bias}&=N^{-1}\sum_i(\widehat y_i-y_i),&
 \mathrm{WAPE}&=\frac{\sum_i|\widehat y_i-y_i|}{\sum_i y_i}.
\end{align}
MAE is the primary criterion; RMSE emphasizes large errors, and negative bias indicates underestimation \citep{hyndman2006}. For intervals, we report observed coverage, mean width and the interval score \citep{gneiting2007},
\begin{equation}
 \mathrm{IS}_\alpha(L,U;y)=(U-L)+\frac{2}{\alpha}(L-y)\ind(y<L)
                  +\frac{2}{\alpha}(y-U)\ind(y>U).
 \label{eq:is}
\end{equation}
Lower score rewards narrow intervals while penalizing missed outcomes. After completion, the excess $y_i-\widehat y_i$ and an out-of-interval flag can identify records for plant review:
\begin{equation}
 r_i^{\mathrm{excess}}=y_i-\widehat y_i,\qquad
 B_i=\ind\{y_i<L_i\;\text{or}\;y_i>U_i\}.
 \label{eq:outputs}
\end{equation}

Pairwise differences use circular bootstrap blocks of three observed origin dates, keeping all machines within a date together. There are 3,000 replicates for the original test and 5,000 for the exploratory comparisons. Alongside marginal 95\% intervals, H1--H3 use approximate Bonferroni-adjusted 98.33\% intervals and the seven temporal contrasts use 99.29\% intervals. These adjustments cover the specified contrasts, not all preceding exploratory choices. Machine and weekly results are descriptive strata with their sample sizes retained.

\section{Results}
\subsection{Production patterns and the baseline comparison}
Electricity use varied widely within each machine, alongside differences in dosed mass and product allocation (Figure~\ref{fig:production}). The scatter of individual records shows why neither an average total consumption nor a single intensity value can fully describe this production setting. Differences between machine-level distributions also reflect differences in what the machines processed. They cannot be read directly as differences in mechanical efficiency.

\fig{02_production_and_energy}{production}{1}{Production and electricity across the three machines. (a--c) Dosed mass and pelleting electricity for all 2,742 positive-energy records, with machine-specific medians annotated. (d--f) Mass distributions, dosing-batch frequencies and product-family mix for all 2,745 source records, including the three zero-energy records. Printed values inside the product bars are record counts. Scatter points represent individual operations, without an implied uniform time spacing.}

In the original comparison on 800 development records, Ridge and histogram gradient boosting gave nearly identical MAE: 51.48 and 51.50 kWh, respectively (Table~\ref{tab:development}). ExtraTrees was close behind, while the multilayer perceptron and the three reference estimators had larger errors. The recent-intensity estimator was less accurate than the training-period intensity reference, despite using completed observations from the validation weeks. Recent information was therefore useful only if incorporated in a suitable way; recency alone did not provide a better baseline.

\begin{table}[!htbp]\centering\small
\caption{Original chronological development comparison on the same 800 records. Configurations were fixed; computational budgets were not identical. The recent-SEC method additionally updates its information during validation.}\label{tab:development}
\begin{tabular}{@{}lrrrr@{}}\toprule
Method & MAE & RMSE & Bias & WAPE (\%)\\\midrule
Ridge & 51.48 & 72.31 & $-3.14$ & 19.88\\
Gradient boosting & 51.50 & 73.46 & $-5.72$ & 19.88\\
ExtraTrees & 52.31 & 75.37 & $-4.89$ & 20.20\\
MLP, three-seed mean & 56.64 & 81.36 & $-1.47$ & 21.87\\
Machine SEC & 59.69 & 83.92 & $-6.59$ & 23.05\\
Recent SEC, last five records & 73.51 & 109.55 & 4.29 & 28.38\\
Machine median & 82.61 & 112.78 & $-14.09$ & 31.89\\\bottomrule
\end{tabular}\par\vspace{4pt}\footnotesize MAE, RMSE and bias are in kWh per record.
\end{table}


Ridge was retained for its accuracy and transparent structure, then fitted on the 1,542 eligible development records. On the subsequent 798-record test, it reduced MAE from 61.29 kWh for the machine-specific intensity reference to 52.45 kWh, an improvement of 8.84 kWh or 14.42\% (Table~\ref{tab:test}). The paired 95\% block-bootstrap interval for the MAE reduction was [6.08, 11.50] kWh. RMSE also decreased, from 92.05 to 80.42 kWh. Removing the five test records with durations exceeding 24 hours left MAE at 52.32 kWh for Ridge and 61.15 kWh for the reference, so those records did not account for the overall advantage.

The improvement occurred on all three machines, but their remaining errors were uneven. Pellet~2 had the lowest Ridge MAE, at 43.36 kWh, compared with 53.36 kWh for Pellet~3 and 66.37 kWh for Pellet~1. Pellet~1 also had a mean prediction error of $-53.85$ kWh, indicating substantial underestimation. The observed-versus-predicted plots make this visible: many of its high-consumption records lie below the equality line (Figure~\ref{fig:testscatter}). Thus, the global improvement coexists with a machine-specific limitation that matters for operational use.

\begin{table}[!htbp]\centering\small
\caption{Archived original test results. Lower MAE for Ridge is observed in each machine; the larger negative bias on Pellet~1 remains visible.}\label{tab:test}
\begin{tabular}{@{}llrrrr@{}}\toprule
Population & Method & $N$ & MAE & RMSE & Bias\\\midrule
All machines & Ridge & 798 & 52.45 & 80.42 & $-21.89$\\
All machines & Machine SEC & 798 & 61.29 & 92.05 & $-24.57$\\\addlinespace
Pellet 1 & Ridge & 200 & 66.37 & 105.14 & $-53.85$\\
Pellet 1 & Machine SEC & 200 & 72.31 & 113.90 & $-64.03$\\
Pellet 2 & Ridge & 333 & 43.36 & 61.26 & $-10.85$\\
Pellet 2 & Machine SEC & 333 & 49.29 & 73.24 & $-8.66$\\
Pellet 3 & Ridge & 265 & 53.36 & 80.10 & $-11.63$\\
Pellet 3 & Machine SEC & 265 & 68.05 & 94.79 & $-14.79$\\\bottomrule
\end{tabular}\par\vspace{4pt}\footnotesize Error metrics are in kWh per record. Global WAPE is 19.11\% for Ridge and 22.34\% for SEC.
\end{table}


\fig{04_original_test_scatter}{testscatter}{1}{Frozen Ridge predictions and errors for every record in the original chronological test. (a--c) Observed versus predicted electricity on identical axes, with equality dashed and MAE, RMSE and bias reported in kWh. (d--f) Signed errors by origin date, with zero dashed and each machine's mean error in red. The 200 Pellet~1, 333 Pellet~2 and 265 Pellet~3 records account for all 798 test observations. The residual axes preserve the full error range.}

\FloatBarrier
\subsection{Temporal models on matched records}
The temporal benchmark asks a narrower question: whether 16 completed same-machine records improve the estimate for the next eligible event. Its history rule leaves 584 shared development targets. On this population, the current-input Ridge control achieved MAE 53.55 kWh, while TFT achieved 53.65 kWh (Table~\ref{tab:gpu}). The paired reduction relative to Ridge was $-0.105$ kWh, with a 95\% interval of [$-2.199$, 2.318] kWh. The interval includes advantages in either direction, and provides no evidence of an improvement from TFT. Adding the full history to Ridge increased MAE to 63.95 kWh; N-HiTS reached 65.28 kWh. These results concern matched event histories and should be interpreted separately from the original 800-record comparison.

\begin{table}[!htbp]\centering\small
\caption{GPU temporal benchmark on 584 identical development records. TFT and N-HiTS are three-seed prediction means; foundation-model weights remain frozen. Units are kWh per record except WAPE. These values must not be compared directly with the original 800-record ranking.}\label{tab:gpu}
\begin{tabular}{@{}lrrrr@{}}\toprule
Method & MAE & RMSE & Bias & WAPE (\%)\\\midrule
Ridge: current & 53.55 & 77.01 & $-4.72$ & 20.37\\
TFT & 53.65 & 79.11 & $-11.23$ & 20.41\\
Ridge: full window & 63.95 & 85.95 & $-1.42$ & 24.32\\
N-HiTS & 65.28 & 92.56 & $-3.05$ & 24.83\\
Chronos-2 + covariates & 85.53 & 120.82 & $-8.11$ & 32.53\\
TimesFM 2.5: energy only & 96.40 & 130.72 & $-9.22$ & 36.67\\
Chronos-2: energy only & 100.02 & 136.47 & $-11.60$ & 38.04\\
TimesFM 2.5 + XReg & 112.12 & 644.61 & $28.03$ & 42.65\\\bottomrule
\end{tabular}\end{table}


\fig{08_gpu_comparison}{gpucomparison}{1}{Accuracy and paired differences on the 584 shared development targets. (a) MAE with 95\% block-bootstrap intervals. (b) Reduction in MAE relative to current-input Ridge, so positive values favour the alternative. Pale intervals allow for seven comparisons; dark intervals are unadjusted 95\% intervals. The full interval endpoints are retained, including the wide TimesFM XReg interval. Resampling uses blocks of three observed origin dates.}

Chronos-2 benefited from the production covariates compared with its energy-only version, although its MAE of 85.53 kWh remained above the matched Ridge control. The TimesFM XReg variant had both the highest MAE and a particularly large RMSE. Three Pellet~3 records in the first validation week received predictions of approximately 2,753, 9,239 and 12,203 kWh, against observed values of 185.4, 142.4 and 70.9 kWh. In each local 16-record context, dosed mass was almost constant; the current mass lay about 1,100--1,467 historical standard deviations from the local mean. An independent float64 reconstruction closely reproduced the adapter's extreme extrapolation. All three records remain in the reported metrics.

This behavior helps explain the large gap between MAE and RMSE for that variant. It also distinguishes the local covariate adapter from the globally fitted supervised models, which estimate production effects from hundreds of records. The diagnostic concerns the tested short-context XReg configuration. It does not establish how a longer context or a separately selected, stabilized adapter would perform. The machine and week breakdowns in Figure~\ref{fig:gpustrata}, together with the complete predictions in Figure~\ref{fig:gpuall}, show the variation hidden by the overall ranking.

\fig{09_gpu_stratification}{gpustrata}{1}{MAE by machine and validation week for the same eight methods. Values are kWh per record; the labels give the number of targets in each stratum. The colour scale ends at 130 kWh to make differences among the lower errors readable, while printed values preserve every result above that limit. Machine and week panels are two views of the same 584 targets.}

\fig{10_gpu_all_predictions}{gpuall}{1}{All 4,672 predictions from the matched comparison: eight methods evaluated on 584 identical records. Colours identify the three machines and the dashed diagonal marks equality. The main axes share a 0--1,100 kWh range. The TimesFM XReg panel also includes a full-range inset containing all its predictions, including the three extremes. MAE and RMSE use the complete data in every panel.}

\FloatBarrier
\subsection{Interpreting estimates and their uncertainty}
The selected Ridge model allows each estimate to be traced to its inputs. Grouping the encoded columns by variable gives the untruncated prediction
\begin{equation}
 f_i=\widehat b+\sum_{g\in\mathcal G}C_{ig},\qquad
 C_{ig}=\sum_{k\in g}z_{ik}\widehat\beta_k,
 \label{eq:explain}
\end{equation}
where $\mathcal G$ comprises mass, batch count, hour, weekday, machine, product subfamily, presentation and bag/bulk category. This sum reproduces all 798 stored test estimates to within $1.2\times10^{-13}$ kWh; none requires non-negative clipping. Figure~\ref{fig:explanation} distinguishes the fitted coefficients from the contributions they make to actual records.

At fixed values of the other inputs, the mass coefficient corresponds to 13.30 kWh per additional tonne and the batch-count coefficient to 9.52 kWh per additional dosing batch. Their positive signs agree with an increasing energy requirement as the recorded production quantity increases. Mass and batch count are correlated, however, so the two coefficients are conditional associations rather than separately identified physical constants. The categorical effects require similar care because product, presentation and machine allocation are not randomized. In the fitted one-hot representation, the intercept of 288.09 kWh is an algebraic reference and cannot be interpreted as measured standby consumption. Small machine coefficients likewise do not establish equal intrinsic efficiency.

The real operation introduced in Figure~\ref{fig:framework} provides a concrete example. For 18.002 t dosed in five batches on Pellet~2, the model estimates 255.90 kWh and a fixed 95\% interval of [122.94, 388.85] kWh. The recorded outcome is 267.30 kWh, an excess of 11.40 kWh over the estimate and well inside the interval. This record was chosen near the median observed test energy to illustrate the calculation. Its error is not representative of every operation, as the complete residual plots show.

\fig{11_explanation}{explanation}{1}{The frozen Ridge model from coefficients to individual estimates. (a) Numerical coefficients per training standard deviation. (b) All 16 categorical coefficients in the fitted encoding. (c) Grouped contributions for every test record, with medians in red and interquartile ranges in orange; small vertical jitter separates coincident values. (d) Arithmetic decomposition of source row 2080, the example in Figure~\ref{fig:framework}. The reference combines standardized numerical means with all-zero categorical indicators and need not describe a feasible operating condition. Contributions are conditional model terms, not causal effects.}

Two exploratory extensions helped identify where modest improvements might arise (Table~\ref{tab:explore}). Allowing the mass association to vary by product subfamily reduced development MAE by 1.44 kWh, or 2.80\%. A shrunk correction from recent completed residuals reduced it by 1.71 kWh, or 3.33\%. Both retained positive intervals after adjustment for the three exploratory comparisons. By contrast, allowing mass effects to vary only by machine produced a small reduction of 0.22 kWh, with an interval that crossed zero. Figure~\ref{fig:development} shows these comparisons and their machine and week breakdowns. Product structure and stabilized residual updates are therefore candidates for a subsequent evaluation; these exploratory findings did not replace the frozen baseline.

\begin{table}[!htbp]\centering\small
\caption{Exploratory point-model extensions on the same 800 development records. Positive differences favour the extension. The adjusted intervals address these three comparisons only.}\label{tab:explore}
\begin{tabular}{@{}lrrr@{}}\toprule
Variant & MAE & MAE reduction & 98.33\% interval for reduction\\\midrule
H1: mass by machine & 51.26 & 0.217 & [$-0.082$, 0.581]\\
H2: mass by subfamily & 50.04 & 1.442 & [0.552, 2.481]\\
H3: recent residuals & 49.77 & 1.712 & [0.236, 3.156]\\\bottomrule
\end{tabular}\par\vspace{4pt}\footnotesize All quantities are in kWh per record.
\end{table}


\fig{03_development_and_hypotheses}{development}{1}{Original development results and exploratory extensions. (a) MAE of the seven original methods on 800 common records. (b) Paired MAE reductions relative to Ridge for machine-dependent mass effects (H1), product-dependent mass effects (H2) and recent-residual correction (H3), with 95\% and multiplicity-adjusted 98.33\% intervals. (c,d) Reductions by machine and validation week; positive values favour the extension. These strata describe the same development records.}

Uncertainty remained less uniform than the point comparison. The frozen 90\% and 95\% intervals covered 85.96\% and 93.23\% of test records, with mean widths of 183.63 and 258.85 kWh. Their block-bootstrap 95\% coverage intervals were [81.86, 90.24]\% and [90.53, 95.87]\%, respectively, so nominal coverage was still compatible with sampling uncertainty at the global level. Machine-level 95\% coverage was 90.00\% for Pellet~1, 96.40\% for Pellet~2 and 91.70\% for Pellet~3. Of the 54 records outside those limits, 45 were above and nine below them. These are candidates for review, without an independent diagnosis of their cause (Figure~\ref{fig:alltest}).

\fig{05_all_test_records}{alltest}{1}{Every original test record in chronological order within its machine. Black points show observed electricity, coloured points the frozen estimate and individual vertical lines the fixed 95\% intervals. Red circles identify the 54 outcomes outside their intervals. Horizontal positions are record order, so equal spacing does not imply equal elapsed time. All 798 observations and their intervals are retained, without joining adjacent records into a continuous uncertainty band.}

Updating the interval quantiles from the most recent 200 completed residuals increased coverage to 89.47\% and 94.49\%. It also increased mean width by 19.66\% and 18.85\%, while interval score improved by only 1.12\% and 0.23\% (Table~\ref{tab:intervals}). The weekly profiles in Figure~\ref{fig:uncertainty} show that coverage and width moved together. The number of records outside the 95\% interval fell from 54 to 44, but these counts cannot be interpreted as false alarms without plant adjudication. The update was explored after the original test had been examined and provides a calibration diagnostic rather than a second independent validation.

\begin{table}[!htbp]\centering\small
\caption{Fixed and post-hoc updated intervals on the same 798 archived point predictions. The updated method has not been validated on an independent period.}\label{tab:intervals}
\begin{tabular}{@{}llrrr@{}}\toprule
Method & Nominal & Coverage (\%) & Mean width (kWh) & Interval score (kWh)\\\midrule
Fixed & 90\% & 85.96 & 183.63 & 381.79\\
Updated, 200 residuals & 90\% & 89.47 & 219.74 & 377.51\\
Fixed & 95\% & 93.23 & 258.85 & 497.64\\
Updated, 200 residuals & 95\% & 94.49 & 307.64 & 496.48\\\bottomrule
\end{tabular}\end{table}


\fig{06_uncertainty_tradeoff}{uncertainty}{1}{Coverage, width and interval score for fixed calibration and the exploratory rolling update. (a,b) Machine-specific coverage against mean interval width at nominal 90\% and 95\% levels; arrows connect fixed intervals (circles) to updated intervals (triangles). (c,d) Weekly 95\% coverage and interval score, with sample sizes shown and partial weeks at the edges. Dashed horizontal lines indicate nominal coverage. Lower interval score is preferable. Both policies use the same 798 fixed point predictions.}

\section{Discussion}
The clearest result is the improvement obtained by accounting for production conditions. A regularized linear model reduced error against the machine-level intensity reference on all three machines in a later industrial period. This gives the baseline a practical role in comparing completed operations with different masses and product characteristics. The sizeable Pellet~1 underestimation nevertheless shows that a good global result does not remove the need to examine individual machines. The available records do not resolve whether that bias arises from omitted operating conditions, product allocation, changing consumption or measurement boundaries. A baseline fitted to historical operation also reflects any inefficiencies already present in that operation; it does not estimate minimum achievable energy use.

The temporal comparison suggests that, in this dataset, current production characteristics carry much of the usable predictive information. TFT and the current-input Ridge control were close, whereas the full-window Ridge and N-HiTS had larger errors. This pattern is consistent with limited additional value from the tested histories, but it does not identify a causal explanation for an architecture's performance. Each event aggregates a different mass over a different duration, and the event index does not represent a uniform time grid. Moreover, the strict history reset removed 27\% of the original validation targets as warm-up, narrowing the evaluated population. Longer histories, elapsed-time representations and broader architecture selection could change the comparison, but would require their own development protocol. The present fixed configurations do not establish a general ordering of forecasting model families.

The foundation-model diagnostics sharpen this point. Covariate adaptation and normalization influence the final industrial prediction alongside the pretrained weights. TimesFM XReg extrapolated severely when its local mass history was nearly constant; Chronos-2 with covariates showed sensitivity to device-dependent normalization of constant channels (Appendix~\ref{sec:computing}). These observations identify concrete issues to investigate before operational use. They also explain why merely making a model accept the available covariates is insufficient evidence that the complete prediction pipeline is stable for this production setting.

The small exploratory gains from product-dependent mass effects and residual correction offer a more focused next step. Product subfamily can partly distinguish operating demands that a single mass coefficient averages together. Recent residuals may capture persistent departures from the historical relation, provided that updates use only completed outcomes and are sufficiently stabilized. Neither interpretation establishes a causal product effect or equipment change. The favorable results support prespecifying these two candidates for a new period, with separate assessment of machine-level bias, rather than replacing the frozen model on the basis of additional comparisons with an already examined test.

For prospective forecasting, the main unresolved issue is when the inputs become known. The consolidated workbook contains actual dosed mass and batch totals, without plan versions or publication timestamps. Those fields are suitable for retrospective production adjustment, but may only become complete after processing has started. Overlapping records further complicate the relation between a source row and a physical production boundary. Confirming meter scope, stable order identifiers and outcome publication times would therefore add more value to the next evaluation than further refinement of an unverified timestamp proxy. Real time-resolved observations would also be needed to test multiple clock-time horizons, machine inactivity or within-operation transitions. The current one-record task cannot resolve those questions.

Uncertainty should be judged by what the intervals enable in practice. Here, updating recent calibration increased coverage largely through wider bounds and gave little improvement in interval score. The trade-off limits the value of treating every out-of-interval record as an operational signal. A new consecutive industrial period should retain all eligible operations starting within it and follow long operations to completion, with the baseline, adjustment rules and calibration policy fixed beforehand. Such a protocol would provide the evidence needed to integrate the baseline into a digital twin \citep{tao2019}. The interface could retain the input version and model version for each issued estimate, then attach the measured outcome when it becomes available. This would preserve a usable distinction between an earlier forecast and a later production-adjusted assessment, while supporting plant review of unusually high consumption.

\section{Conclusions}\label{sec:conclusions}
A production-adjusted Ridge model gave a more accurate and more traceable baseline for pelleting electricity than the operational references evaluated here. In the frozen chronological test it reduced MAE by 14.42\% relative to line-specific SEC while providing a reproducible additive decomposition of every estimate. The improvement concerns the estimation of expected electricity under the recorded production conditions; it is not a reduction in plant energy use or evidence of achievable savings.

Under the tested 16-record, same-line, within-week representation, completed operating histories did not improve on current-record covariates, and neither TFT, N-HiTS nor the tested TimesFM~2.5 and Chronos-2 configurations beat the current-input control. Production adjustment helped on all three lines, most consistently on Pellet~2 and Pellet~3; Pellet~1 improved on SEC, but its systematic underestimation calls for line-specific recalibration or additional predictors before prospective use.

The evidence supports retrospective assessment at the closure of recorded operations, where the baseline provides an expected value, an empirically calibrated residual interval and a traceable deviation for identifying candidates for technical review. Prospective supervisory use requires validation on a new consecutive production period with prespecified line-level criteria and verified record identity, measurement boundaries, predictor availability and calibration rules. Until then, automatic control, modification of process set-points and operational digital-twin integration remain outside the demonstrated scope of the study.

\section*{CRediT authorship contribution statement}
%%% DRAFT roles, to be confirmed by all authors before submission.
\textbf{Vladimir Sousa Santos:} Conceptualization, Methodology, Investigation, Writing -- original draft, Writing -- review \& editing, Project administration, Funding acquisition.
\textbf{Juan Jos\'e Cabello Eras:} Conceptualization, Resources, Supervision, Writing -- review \& editing.
\textbf{Eder El\'{\i}as Herrera Rodi\~no:} Investigation, Data curation, Resources, Writing -- review \& editing.
\textbf{Francisco Manuel Arrabal-Campos:} Methodology, Software, Formal analysis, Data curation, Validation, Visualization, Writing -- original draft, Writing -- review \& editing.

\section*{Funding}
This work was carried out within the postdoctoral research stay of V. Sousa Santos at the University of Almer\'{\i}a, supported by the Asociaci\'on Universitaria Iberoamericana de Postgrado (AUIP) mobility programme (2026).

\section*{Declaration of competing interest}
The authors declare that they have no known competing financial interests or personal relationships that could have appeared to influence the work reported in this paper.

\section*{Acknowledgements}
The authors thank Alimentos Finca S.A.S. (Ci\'enaga de Oro, C\'ordoba, Colombia) for providing the production and electricity records of its pelleting lines.

\section*{Data availability}
The code, cleaned records, frozen model, archived predictions, model configurations, figure data and LaTeX sources needed to reproduce every table and figure are openly available in the GitHub repository \url{https://github.com/fmarrabal/pelleting-energy-baseline-real-data}. The versioned research archive attached to its releases also contains the original plant workbook and the complete experimental outputs, linked to the manuscript by retained SHA-256 fingerprints. The industrial records were provided by the plant for research purposes; their reuse beyond reproduction of this study requires agreement with the data owner. Pretrained checkpoints are identified by repository and revision and are not redistributed.

\appendix
\section{Computational and numerical reproducibility}\label{sec:computing}
The GPU experiment used Python 3.13.13, PyTorch 2.14.0 with CUDA 13.0, NeuralForecast 3.2.2, the TimesFM package 3.0.2, Chronos-forecasting 2.3.2 and scikit-learn 1.9.0. The pretrained checkpoint revisions were \texttt{1d952420fba87f3c6dee4f240de0f1a0fbc790e3} for TimesFM 2.5 and \texttt{29ec3766d36d6f73f0696f85560a422f50e8498c} for Chronos-2. TimesFM received 16 observed context records, padded to a maximum context of 32, with a configured maximum horizon of 32; only its first forecast position was evaluated. Chronos used float32, disabled cross-query learning and returned the median forecast. Model configurations, checkpoints and record-level outputs are retained in the repository; pretrained weights are identified by revision rather than redistributed.

All 128,896 logged donor relationships passed the checks for same-pelleting-line membership, partition membership and completion before the target origin; the total includes repeated relationships across nested and expanding blocks. Source-data, frozen-model and final-prediction fingerprints remained unchanged, and the shared CPU and GPU data, window and encoding functions were checked for consistency. These checks concern the recorded completion-time proxy, since actual publication delays are not available in the source.

\fig{12_numerical_and_seeds}{numerics}{1}{Numerical and training variability. (a) Every paired CPU/GPU Chronos-2 covariate output as a signed difference against the CPU prediction. (b) All three GPU seeds (11, 29 and 47) and their prediction-mean ensemble for TFT and N-HiTS, with the matched Ridge reference; averaging predictions is not the same as averaging the three MAEs. These diagnostics accompany the primary comparison without selecting a preferred seed or device.}

Within the GPU run, the 24 saved neural checkpoints were reloaded and compared with their stored inference, together with four Chronos covariate batch-size checks. All 28 checks passed the recorded tolerances with a maximum difference of $7.63\times10^{-5}$ kWh. Fresh CPU and GPU neural fits were less closely aligned, as optimization can follow different trajectories even with the same integer seed: TFT's CPU and GPU MAEs were 53.25 and 53.65 kWh, and neither established an advantage over the matched Ridge control under the paired intervals. Figure~\ref{fig:numerics} shows all three GPU seeds and their ensemble.

For frozen Chronos-2 with covariates, the median absolute CPU--GPU prediction difference was 0.86 kWh and the maximum 57.04 kWh, and its overall MAE moved from 85.65 to 85.53 kWh. Maximum prediction differences for both frozen TimesFM variants and for univariate Chronos stayed below $1.9\times10^{-4}$ kWh. To trace the covariate sensitivity, the installed Chronos preprocessing and instance normalization were applied to the same 584 contexts, giving 6,424 target/covariate channels. Among the 1,097 constant historical channels, 694 showed a CPU--GPU difference above 0.1 in normalized context values, whereas no non-constant channel exceeded that threshold; the largest difference was 0.88137 in the transformed coordinates.

The implementation substitutes a small scale only when the calculated scale is exactly zero, so tiny round-off variances in constant channels can be amplified before the inverse-hyperbolic-sine transform. This mechanism is consistent with the observed sensitivity, although it is not an end-to-end causal demonstration for each prediction. The reported benchmark keeps the native preprocessing and original outputs; a corrective normalization policy would require a separately selected and evaluated experiment.

\FloatBarrier
\section*{Declaration of generative AI and AI-assisted technologies in the writing process}
During the preparation of this work the authors used Trinka as a language-editing assistant to improve the readability and English of the manuscript, and Claude Code as a coding assistant for writing and documenting the analysis scripts. After using these tools, the authors reviewed and edited the content as needed and take full responsibility for the content of the published article.


\bibliographystyle{elsarticle-num}
\bibliography{references}
\end{document}
